\subsection{Dinâmica de Atitude}
Utilizando as equações de Euler aplicadas a um corpo rígido, é possível descrever a atitude de uma espaçonave. 
Considerando que o movimento é dado no centro de massa do corpo e sendo expresso no sistema da espaçonave, as equações de movimento rotacional assumem a forma:

\begin{equation}
\mathbf{I}\,\dot{\vec{\omega}}
+ \vec{\omega}
\times \big(\mathbf{I}\,\vec{\omega}\big)
= \vec{M},
\label{eq:euler_corpo_rigido}
\end{equation}

em que $\mathbf{I}$ é a matriz de inércia em relação ao centro de massa, $\vec{\omega}$ é o vetor de velocidade angular do corpo e $\vec{M}$ é o vetor dos momentos aplicados (torques externos).

\subsubsection{Matriz de inércia}
A matriz de inércia equivalente da espaçonave não é constante, porque a movimentação do MR altera a distribuição de massa.
Sendo assim, a matriz total pode ser expressa como:
\begin{equation}
\mathbf{I}(t)
= \mathbf{I}_{base}
+ \mathbf{I}_{manip(t)},
\label{eq:inercia_total}
\end{equation}

em que $\mathbf{I}_{base}$ é a matriz associada à base da espaçonave e $\mathbf{I}_{manip(t)}$ é a contribuição variável do manipulador.
Os torques de controle e os torques de reação são inclusos nos momentos aplicados sob o corpo.

\subsubsection{Contribuição do manipulador robótico acoplado à base}
Ao movimentar-se, o MR produz variações na inércia do sistema a dinâmica acoplada pode ser representada por:

\begin{equation}
\mathbf{I}(t)\,\dot{\vec{\omega}}
+ \vec{\omega}
\times \big(\mathbf{I}(t)\,\vec{\omega}\big) 
=
\vec{M}_{controle}
+ \vec{M}_{manip},
\label{eq:dinamica_acoplada}
\end{equation}

onde $\vec{M}_{controle}$ são os torques de controle aplicados à espaçonave e $\vec{M}_{manip}$ são os torques de reação gerados pelas juntas do manipulador.  
% Portanto, a modelagem da dinâmica de atitude demanda considerar de forma simultânea: a inércia variável, os torques de controle e os efeitos dinâmicos do manipulador robótico acoplado à base.